\section{Hipóteses}
\label{sec:hipoteses}

A evidência de cada hipótese está em \texttt{gerado/pi5/hypotheses.md}; o
veredicto é dos autores.

\begin{table}[ht]
  \centering
  \caption{Veredicto das hipóteses da Parte~1.}
  \label{tab:hipoteses}
  \footnotesize
  \begin{tabularx}{\textwidth}{@{}L{5.6cm}YL{2.2cm}@{}}
    \toprule
    \textbf{Hipótese} & \textbf{Evidência} & \textbf{Veredicto} \\
    \midrule
    H1. O envelope AES-256-GCM tem um custo por mensagem pequeno face à latência
      da rede & K1$\to$K3: +7,2~$\mu$s de processador (IC $[6{,}8; 7{,}8]$) e
      +25~$\mu$s no p50 (149~$\mu$s em K1); débito máximo de 40\,000 para
      30\,000~mensagens/s & Confirmada, com ressalva \\
    H2. Validade, anti-repetição e chave--tópico têm custo marginal quase nulo &
      K3$\to$K4 +0,6~$\mu$s e K3$\to$K5 +0,2~$\mu$s (IC incluem 0); K6$\to$K7
      +1,3~$\mu$s (3\,\%); p99 sem diferença mensurável & Confirmada \\
    H3. Ed25519 custa claramente mais do que AES-GCM por mensagem &
      K8$\to$K9: +151,5~$\mu$s (4,7 vezes), +464~$\mu$s no p99, débito máximo seis
      vezes menor (5\,000 contra 30\,000~mensagens/s) & Confirmada \\
    H4. O custo do TLS concentra-se no estabelecimento da ligação &
      Por ligação: 4,2 contra 0,8~ms de processador. Por mensagem: K0$\to$K1
      +7,9~$\mu$s e K10$\to$K7 +5,5~$\mu$s; débito máximo de 60\,000 para
      40\,000~mensagens/s & Refutada \\
    H5. Sob inundação, a limitação antes da criptografia reduz o processador do
      \emph{gateway} & Inundação a 300~msg/s: K7 5,2\,\%, K8 5,1\,\% de
      processador (40,5 e 39,7~$\mu$s por mensagem); K8 descarta 1\,559
      mensagens antes da criptografia & Não confirmada \\
    \bottomrule
  \end{tabularx}
\end{table}

\paragraph{H1 (confirmada, com ressalva).} Em valor absoluto, o envelope custa
pouco: 7,2~$\mu$s de processador e 25~$\mu$s de latência mediana, mil vezes menos
do que $L_{\max}$ e menos do que a latência de qualquer rede sem fios típica da
IoT. Mas, nesta rede local cablada, a latência é tão baixa que o envelope
representa 17\,\% do p50 de K1, e reduz o débito máximo em 25\,\%. A hipótese
vale para a latência; para o débito, o envelope não é desprezável.

\paragraph{H4 (refutada).} A negociação TLS com autenticação mútua custa
3,4~ms de processador a mais por ligação, mas o TLS custa também 7,9~$\mu$s por
mensagem, tanto como o envelope inteiro (7,2~$\mu$s), e é o que mais reduz o
débito máximo (de 60\,000 para 40\,000~mensagens/s). Ao fim de cerca de 430
mensagens, o custo por mensagem ultrapassa o da negociação; a 1\,000~mensagens/s,
isso acontece em menos de meio segundo, e um sensor a 10~Hz chega lá em 43~s. Em
ligações duradouras, que são as do MQTT, o custo do TLS está nas mensagens, não
no estabelecimento. O custo por mensagem é menor com o envelope ativo (5,5 contra
7,9~$\mu$s, Secção~\ref{sec:custo}), mas continua a dominar.

\paragraph{H5 (não confirmada).} A limitação funciona como previsto: descarta o
excesso antes de qualquer verificação criptográfica (1\,559 de 1\,814 mensagens
da inundação). Mas o processador do \emph{gateway} é o mesmo em K7 e K8 (5,2 e
5,1\,\%). Com AES em \emph{hardware}, verificar uma mensagem custa poucos
microssegundos: o que a limitação poupa, cerca de 260 verificações por segundo,
é da ordem de 0,2\,\% de um núcleo, abaixo do ruído. Com verificações mais caras,
como as assinaturas, a poupança poderia ser maior, mas a matriz não tem uma
configuração com assinaturas e sem limitação que o meça. A hipótese, tal como formulada para o
envelope AES-GCM, não se confirma. O que a limitação mostrou foi a proteção da
aplicação: sem ela (K7), as mensagens da inundação são todas entregues; com ela
(K8), o excesso é descartado (Secção~\ref{sec:protecao}).
